\section{本章小结}
\label{sec:ls-inner-product-spaces-summary}

本章回答了章首的问题：内积在一般线性空间上规定了角度，由此可以判断垂直、求出最好的近似，并选出让坐标计算最简单的基。

内积给出长度、正交与夹角。Cauchy--Schwarz不等式保证夹角有定义；内积诱导的长度满足三角不等式，因而是范数，并给出距离。并非每个范数都来自内积，平行四边形恒等式可以检验这一点；有限维内积空间是Hilbert空间。

Gram矩阵记录一组向量的两两内积。它总是半正定的，并且恰在这组向量线性无关时正定；反过来，每个半正定矩阵都是某组向量的Gram矩阵，这又推出逐元素乘积保持半正定。基的Gram矩阵把一般内积写成坐标的计算公式，换基时按合同变换改变。

正交分解把向量拆成子空间中的部分与垂直的部分。标准正交基使坐标可以由内积直接读出，Parseval恒等式说明长度的平方等于坐标的平方和；正交矩阵是标准正交基之间的换基矩阵，求逆只需转置，并保持欧氏长度。Gram--Schmidt正交化保证标准正交基总存在，QR分解把这一过程写成矩阵形式。正交投影是子空间中唯一的最近点；矩阵是正交投影，当且仅当它幂等且对称。

这些结论为后续各章提供几何工具。第~\ref{chap:dual-spaces}章用内积把线性泛函表示为向量、把对偶映射改写为伴随；第~\ref{chap:eigen-spectral-decomposition}章在正交换基下寻找线性作用的简单表示，正交矩阵与自伴算子在那里尤为重要；最小二乘、QR求解与正交投影则贯穿线性方程各章。
